\section{The exact Euler domains and their receiving maps}


This independent addendum checks (PM20)--(PM23) in
\path{PRIME_MEMORY_AND_ACTUAL_THETA_RECEIVER.tex}.
It uses the complete module and map calculations in
\path{moment_first_coinvariants_verification.tex},
(MC5)--(MC7), (MC21)--(MC28), and the marked Hilbert receiver
constructed in \path{connes_marked_quotient_verification.tex},
(MQ15)--(MQ18). It specifies the logarithmic Fourier convention
and the full orthogonal-sum domain explicitly. The differential
operator, original real coordinate, moments, and retained constant
$1/4$ are unchanged.

\subsection{The exact real inverse and the full adelic quotient}

Write $\mathscr S=\mathcal S(\mathbb R)^{\mathrm{even}}$,
$m(h)=(h(0),\int h)$, $\mathscr S_0=\ker m$,
$E=x\partial_x$, and $\Delta=E(1+E)$. The letter $E$ in
this addendum denotes the differential operator, not periodization.
First $\Delta\mathscr S\subset\mathscr S_0$: its evaluation at
zero is zero and integration by parts gives
$\int\Delta g=-\int(1+E)g=0$.

\begin{theorem}[The complete inverse]
For $f\in\mathscr S_0$, put for $x>0$
\[
 u(x)=-\int_x^\infty f(t)t^{-1}\,dt,
 \qquad g(x)=-x^{-1}\int_x^\infty u(t)\,dt.
 \tag{LD1}\label{ld:eq:inverse}
\]
Their even extensions are Schwartz functions, and $\Delta g=f$.
This defines the inverse of the continuous linear isomorphism
$\Delta:\mathscr S\to\mathscr S_0$.
\end{theorem}
\begin{proof}
Because $f$ is even and $f(0)=0$, its quotient $f(t)/t$ extends
smoothly and oddly across zero and vanishes there. Differentiating
the first integral gives $u'=f/x$, hence $Eu=f$. Its even extension
is smooth. Its tail and each derivative decay rapidly at infinity.
Fubini is valid since the corresponding double absolute integral
equals $\int_0^\infty|f(t)|dt$, and gives
$\int_0^\infty u=-\int_0^\infty f=0$. Therefore near zero
\[
 g(x)=x^{-1}\int_0^xu(t)\,dt=\int_0^1u(tx)\,dt.
\]
This proves smooth even extension of $g$, including all derivatives.
The tail expression in \eqref{ld:eq:inverse} proves rapid decay of
it and its derivatives. Differentiating $xg$ gives
$(1+E)g=u$ and hence $\Delta g=f$.

The kernel of $E$ in Schwartz space is zero: a solution is constant
on each half-line and rapid decay forces each constant to vanish.
The kernel of $1+E$ is zero too, because $xg$ is constant on each
half-line. Their product therefore has zero kernel on Schwartz
space. This proves the inverse assertion.
For continuity of the inverse, on $0\le t\le1$ use
$|f(t)|\le t^2\sup_{[0,1]}|f''|/2$. The integral defining
$u(0)$ and the derivatives of $u$ are then bounded by finitely
many Schwartz seminorms of $f$. The displayed average for $g$
does the same near zero. On $x\ge1$ the tail formulas and their
derivatives give any requested weighted bound from finitely many
weighted bounds for $f$ and its derivatives. Thus every Schwartz
seminorm of $g$ is controlled by finitely many of those of $f$.
The direct differential operator is continuous, proving the
claimed topological isomorphism as well as its explicit inverse.
\end{proof}

In the original spherical module
$V^{K,\mathrm{even}}=R\otimes\mathscr S$, rational scaling acts
on $R$ and $E$ acts on the second factor: Euler differentiation
commutes with $x\mapsto x/a$. The entire range is therefore
$R\otimes\mathscr S_0$. In the original moment-first module
the exact quotient is
\[
 M_0^{\mathrm{even}}/\Delta(V^{K,\mathrm{even}})
     \xrightarrow{\sim}I\otimes\C^2,
 \qquad\sum_at_a\otimes h_a\longmapsto
                        \sum_at_a\otimes m(h_a).
 \tag{LD2}\label{ld:eq:adelic-cokernel}
\]
The zero total moment is exactly the augmentation-zero condition.
Its kernel and surjectivity follow coefficient by coefficient from
\eqref{ld:eq:inverse} and the two explicit moment lifts in (MC8).
This verifies (PM21) on the full original coefficient module.

The domain of the induced operator on moment-first coinvariants
is different from the domain in \eqref{ld:eq:adelic-cokernel}.
Put $P=\Lambda\otimes\C^2$. The canonical splitting (MC5)--(MC7)
gives the exact operator
\[
 (M_0)_\Gamma=\mathscr S_0\oplus P,
 \qquad \Delta=(\Delta|_{\mathscr S_0})\oplus0.
 \tag{LD3}\label{ld:eq:coinvariant-operator}
\]
Indeed the unramified section commutes with the operator. On a
prime class, applying $\Delta$ gives
$[(\alpha_p-1)s(\Delta h)]=0$, because $\Delta h\in\mathscr S_0$
and this is an actual relation imposed after the moments.
Equivalently $E$ acts by $0$ and $-1$ on the two pole coordinates,
so $E(1+E)$ kills both. The same argument kills every higher pole
homology group (MC11).

Consequently its exact kernel and cokernel are
\[
 \ker\Delta=P,\qquad
 \operatorname{coker}\Delta\simeq\C^2\oplus P.
 \tag{LD4}\label{ld:eq:coinvariant-defects}
\]
For the second equality the actual quotient map on the first
factor is $f\mapsto m(\Delta^{-1}f)$ from $\mathscr S_0$ to
$\C^2$. It is onto by applying $\Delta$ to arbitrary moment
lifts, and its kernel is exactly $\Delta\mathscr S_0$.
This proves every map in \eqref{ld:eq:coinvariant-defects}.

\subsection{The Haar realization, constants, and full domain}

On a Schwartz coefficient in additive real $L^2$, integration
by parts gives the exact formulas
\[
 E^*=-1-E,\qquad
 \Delta=-E^*E=(E+\tfrac12)^2-\tfrac14,
 \qquad
 \langle f,\Delta f\rangle
     =-\|(E+\tfrac12)f\|^2-\tfrac14\|f\|^2.
 \tag{LD5}\label{ld:eq:haar-sign}
\]
The boundary term is $x\overline f g$, which vanishes at zero
and infinity. The middle identity uses the skew adjointness of
$E+1/2$ on this core. Every orthogonal finite-place projection
and every marked arithmetic mean acts on another factor and
commutes with these operations. Thus the identities hold on the
actual zero-sum marked arithmetic channels, including rational
group labels.

On the positive and negative half-lines define respectively
\[
 (U_+f)(t)=e^{t/2}f(e^t),\qquad
 (U_-f)(t)=e^{t/2}f(-e^t).
 \tag{LD6}\label{ld:eq:log-map}
\]
Changing variables gives unitary maps to $L^2(\mathbb R,dt)$.
Direct differentiation sends $E+1/2$ to $\partial_t$.
For the multiplier $-\xi^2-1/4$ the required logarithmic Fourier
convention is explicitly
\[
 (\mathcal F_{\mathrm{ang}}g)(\xi)
    =(2\pi)^{-1/2}\int_{\mathbb R}e^{-it\xi}g(t)\,dt.
 \tag{LD7}\label{ld:eq:angular-fourier}
\]
It sends $\partial_t$ to $i\xi$. If instead the earlier real
Fourier convention $\int e^{-2\pi it\xi}g(t)dt$ is retained,
the multiplier is $-4\pi^2\xi^2-1/4$. These are the exact
coordinate change $\xi_{\mathrm{ang}}=2\pi\xi$, not different
operators or a change of the retained constant $1/4$.

Let $\mathscr K_a$ include all finite-place and zero-sum marked
coefficient factors in an actual arithmetic channel. The full
Hilbert sum, after \eqref{ld:eq:log-map}--\eqref{ld:eq:angular-fourier},
has the selfadjoint realization
\[
 \begin{gathered}
 (\Delta h)_{a,\pm}(\xi)=(-\xi^2-\tfrac14)h_{a,\pm}(\xi),\\
 \Dom\Delta=\left\{h:
   \sum_{a,\pm}\int_{\mathbb R}
    (\xi^2+\tfrac14)^2\|h_{a,\pm}(\xi)\|_{\mathscr K_a}^2
           \,d\xi<\infty\right\}.
 \end{gathered}
 \tag{LD8}\label{ld:eq:full-domain}
\]
Here $h$ also belongs to the Hilbert sum itself. For an even
real sector the unitary map from the entire even real $L^2$
space to one copy is $f\mapsto\sqrt2\,e^{t/2}f(e^t)$;
the same multiplier
and weighted domain hold. The summability in
\eqref{ld:eq:full-domain} is essential: membership in each individual
channel's domain does not suffice for an infinite orthogonal sum.

To prove selfadjointness, test the adjoint against compactly
supported frequency functions with finitely many coefficient
channels. The test determines its value to be the same real
multiplier and imposes exactly \eqref{ld:eq:full-domain}.
Conversely every vector in that domain satisfies the adjoint
identity by Cauchy--Schwarz. Its nonreal resolvents are bounded
multipliers $(-\xi^2-1/4-z)^{-1}$, which verifies the entire
resolvent assertion as well. Finite channel truncation and
frequency cutoff approximate every domain vector in graph norm.
Smooth compact frequency approximations then give Schwartz
functions in $t$. Multiplying those by smooth cutoffs tending
to one approximates them with their first two derivatives in
$L^2(dt)$. Their inverse images under \eqref{ld:eq:log-map} are
smooth and compactly supported away from $x=0$. Thus the original
finite Schwartz-channel operator is a core for this realization.

On every nonzero channel, frequencies in arbitrarily small
neighbourhoods of a prescribed $\xi_0$ give approximate
eigenvectors of value $-\xi_0^2-1/4$. Outside
$(-\infty,-1/4]$ the reciprocal multiplier is bounded.
Every level set of the multiplier has measure zero, so no
nonzero Hilbert eigenvector exists. This proves the full spectrum
$(-\infty,-1/4]$ and all domain claims in the stated angular
convention. In particular zero is in its resolvent set.

\subsection{The exact comparison with the retained prime kernel}

The operator identities above do not identify their different
receiving modules. Their actual maps, already proved in (MC21)
and (MC22), give the diagram
\[
 \begin{tikzcd}[column sep=large]
 (M_0)_\Gamma=\mathscr S_0\oplus P
 \arrow[r,"\delta"]\arrow[rr,bend right=20,"0"']
 &D_{\Gamma,0}\arrow[r]
 &\overline D^{\,K}/\overline{\delta V^K}.
 \end{tikzcd}
 \tag{LD9}\label{ld:eq:zero-composite}
\]
The first map has kernel exactly $\mathscr S_0$ and image
exactly $P$. The second map kills exactly $P$ on the algebraic
marked quotient, and its image is the faithfully embedded dense
arithmetic mixed-place quotient. Therefore the displayed composite
is exactly zero. All arrows commute with $E$ and $\Delta$ on
their stated algebraic domains. On the prime module $\Delta=0$;
on the attained positive Haar quotient its closed realization
has the strictly negative spectral gap just proved. This
agreement is accounted for by the exact kernel in
\eqref{ld:eq:zero-composite}. It supplies no identification of the
theta receiver with the Haar arithmetic receiver and no transfer
of a zeta-zero assertion between them.

